\chapter{تعامد}
\label{ch:orthogonality}

در فصل قبل با ضرب داخلی «عمود بودن» را در هر فضایی تعریف کردیم. در این فصل می‌بینیم که بردارهای دوبه‌دو عمود چقدر کار را ساده می‌کنند: خودبه‌خود مستقل‌اند و ضرایب ترکیب خطی با یک ضرب داخلی به دست می‌آیند. بعد زیرفضاهای متعامد و مکمل متعامد را می‌شناسیم و با الگوریتم \textbf{گرام-اشمیت} از هر پایه، یک پایه‌ی یکامتعامد می‌سازیم.

\section{مجموعه‌های متعامد}

\begin{definition}[متعامد و یکامتعامد]
بردارهای $\va_1,\dots,\va_k$:
\begin{itemize}
  \item \term{متعامد}{Orthogonal} (دوبه‌دو عمود) هستند اگر $\va_i\perp\va_j$ برای هر $i\neq j$؛
  \item \term{نرمال‌شده}{Normalized} هستند اگر $\norm{\va_i}=1$ برای هر $i$؛
  \item \term{یکامتعامد}{Orthonormal} هستند اگر هر دو شرط برقرار باشد؛ یعنی
  \[
  \va_i\T\va_j=\delta_{ij}=\begin{cases}1 & i=j\\ 0 & i\neq j\end{cases}.
  \]
\end{itemize}
پایه‌ی $\mathcal{B}$ از فضای ضرب داخلی $V$ یک \term{پایه‌ی یکامتعامد}{Orthonormal Basis} است اگر برای $\vv\neq\vw$ در $\mathcal{B}$، $\inner{\vv}{\vw}=0$ و برای هر $\vv\in\mathcal{B}$، $\norm{\vv}=1$.
\end{definition}

\begin{intuition}[عمود یعنی چه؟]
$\vu$ بر $\vv$ عمود است اگر فاصله‌اش تا $\vv$ و تا $-\vv$ برابر باشد. با بسط دادن،
$\norm{\vu-\vv}^2-\norm{\vu+\vv}^2=-4\vu\T\vv$؛
پس این دو فاصله برابرند اگر و تنها اگر $\vu\T\vv=0$.
\begin{center}
\begin{tikzpicture}[scale=0.9,>=Latex]
  \draw[thick, primary] (-2.8,0)--(2.8,0);
  \draw[thick, primary] (0,-0.3)--(0,2.6);
  \fill (0,2.3) circle (2pt) node[above right] {$\vu$};
  \fill (-2,0) circle (2pt) node[below] {$\vv$};
  \fill (2,0) circle (2pt) node[below] {$-\vv$};
  \fill (0,0) circle (2pt) node[below right] {$\zero$};
  \draw[excolor] (0,2.3)--(-2,0) node[midway, above left] {\small $\norm{\vu-\vv}$};
  \draw[excolor] (0,2.3)--(2,0) node[midway, above right] {\small $\norm{\vu-(-\vv)}$};
  \draw (0,0.25)--(0.25,0.25)--(0.25,0);
\end{tikzpicture}
\end{center}
\end{intuition}

در هر فضای ضرب داخلی، $\vv$ و $\vw$ متعامدند اگر $\inner{\vv}{\vw}=0$. \textbf{قضیه‌ی فیثاغورس} (قضیه‌ی \ref{thm:norm-sum}) می‌گوید:
\begin{keyformula}
\centering
$\vu\perp\vv\iff\norm{\vu+\vv}^2=\norm{\vu}^2+\norm{\vv}^2$.
\end{keyformula}

\begin{theorem}\label{thm:orth-indep}
اگر $S=\{\va_1,\dots,\va_k\}$ مجموعه‌ای متعامد از بردارهای \textbf{ناصفر} در $\R^n$ باشد، $S$ مستقل خطی است و پایه‌ای برای $\Span(S)$ است. به‌ویژه اگر $k=n$، $S$ پایه‌ای برای $\R^n$ است.
\end{theorem}
\begin{proof}
فرض کنید $c_1\va_1+\dots+c_k\va_k=\zero$. دو طرف را در $\va_j$ ضرب داخلی می‌کنیم. چون $\va_i\T\va_j=0$ برای $i\neq j$، از جمع فقط یک جمله می‌ماند:
\[
0=\va_j\T(c_1\va_1+\dots+c_k\va_k)=c_j\,\va_j\T\va_j=c_j\norm{\va_j}^2 .
\]
چون $\va_j\neq\zero$، $\norm{\va_j}^2>0$ و $c_j=0$. پس $S$ مستقل است و چون $\Span(S)$ را تولید می‌کند، پایه‌ی آن است. اگر $k=n$، $S$ شامل $n$ بردار مستقل در $\R^n$ است و طبق قضیه‌ی \ref{thm:basis-props} پایه است.
\end{proof}

\begin{note}
پس هر مجموعه‌ی یکامتعامد مستقل است و طبق قضیه‌ی \ref{thm:more-than-n} داریم $k\le n$. اگر $k=n$، بردارها یک \textbf{پایه‌ی یکامتعامد} $\R^n$ هستند.
\end{note}

\begin{corollary}[ضرایب با ضرب داخلی]\label{cor:orth-coeffs}
اگر $\va_1,\dots,\va_k$ متعامد و ناصفر باشند و $\vy=c_1\va_1+\dots+c_k\va_k$، آن‌گاه
\[
c_j=\frac{\vy\T\va_j}{\va_j\T\va_j}.
\]
اگر یکامتعامد باشند، $c_j=\va_j\T\vy$؛ پس $\vy$ ترکیب خطی آن‌هاست اگر و تنها اگر
\begin{keyformula}
\[
\vy=(\va_1\T\vy)\,\va_1+\dots+(\va_k\T\vy)\,\va_k .
\]
\end{keyformula}
\end{corollary}
\begin{proof}
مثل اثبات قبل، $\vy$ را در $\va_j$ ضرب داخلی می‌کنیم: $\vy\T\va_j=c_j\,\va_j\T\va_j$. برای بخش آخر: اگر $\vy$ ترکیب خطی باشد، ضرایبش همین‌هاست؛ و برعکس، اگر تساوی برقرار باشد، $\vy$ ترکیب خطی است.
\end{proof}

دیگر لازم نیست دستگاه حل کنیم؛ هر ضریب فقط یک ضرب داخلی است.

\begin{example}
$\vx=(1,2,3)$ را بر حسب
$\va_1=(0,0,-1)$، $\va_2=\frac1{\sqrt2}(1,1,0)$ و $\va_3=\frac1{\sqrt2}(1,-1,0)$
بنویسید.
\end{example}
\begin{solution}
این سه بردار یکامتعامدند (ضرب داخلی هر دوتا صفر و طول هر کدام ۱ است). پس
$\va_1\T\vx=-3$، $\va_2\T\vx=\frac{3}{\sqrt2}$ و $\va_3\T\vx=-\frac{1}{\sqrt2}$؛ در نتیجه $\vx=-3\va_1+\frac{3}{\sqrt2}\va_2-\frac{1}{\sqrt2}\va_3$. بررسی: $(0,0,3)+(\frac32,\frac32,0)+(-\frac12,\frac12,0)=(1,2,3)$.
\end{solution}

\section{زیرفضاهای متعامد}

\begin{definition}[زیرفضاهای متعامد]
زیرفضاهای $W_1$ و $W_2$ از $V$ \textbf{متعامد} هستند و می‌نویسیم $W_1\perp W_2$، اگر هر بردار $W_1$ بر هر بردار $W_2$ عمود باشد:
$\inner{\vw_1}{\vw_2}=0$ برای همه‌ی $\vw_1\in W_1$ و $\vw_2\in W_2$.
\end{definition}

\begin{theorem}\label{thm:orth-basis-subspaces}
اگر هر بردار پایه‌ی $W_1$ بر هر بردار پایه‌ی $W_2$ عمود باشد، $W_1\perp W_2$.
\end{theorem}
\begin{proof}
اگر $\vw_1=\sum_i a_i\vu_i$ و $\vw_2=\sum_j b_j\vv_j$ (بر حسب دو پایه)، با خطی بودن ضرب داخلی در هر دو آرگومان
$\inner{\vw_1}{\vw_2}=\sum_i\sum_j a_i\conj{b_j}\inner{\vu_i}{\vv_j}=0$.
\end{proof}

پس برای بررسی تعامد دو زیرفضا کافی است بردارهای پایه (یا هر مجموعه‌ی مولد) را بررسی کنیم.

\section{مکمل متعامد}

\begin{definition}[مکمل متعامد]
بردار $\vz$ بر زیرفضای $W$ از $\R^n$ \textbf{عمود} است اگر بر هر بردار $W$ عمود باشد. مجموعه‌ی همه‌ی این بردارها را \term{مکمل متعامد}{Orthogonal Complement} $W$ می‌نامیم:
\[
W^\perp=\{\vz\in\R^n\ :\ \vz\T\vw=0\ \text{\fa{برای هر}}\ \vw\in W\}.
\]
\end{definition}

\begin{example}
اگر $W$ صفحه‌ای گذرنده از مبدأ در $\R^3$ و $L$ خط گذرنده از مبدأ و عمود بر آن باشد، $L=W^\perp$ و $W=L^\perp$.
\begin{center}
\begin{tikzpicture}[scale=0.85,>=Latex]
  \fill[thmcolor!25] (-2.6,-1.1)--(1.4,-1.1)--(2.6,1.1)--(-1.4,1.1)--cycle;
  \draw[thmcolor] (-2.6,-1.1)--(1.4,-1.1)--(2.6,1.1)--(-1.4,1.1)--cycle;
  \node[thmcolor!70!black] at (2.2,-0.8) {$W$};
  \draw[very thick, primary] (0,0)--(0,2.4) node[right] {$L=W^\perp$};
  \draw[very thick, primary, dashed] (0,0)--(0,-1.1);
  \draw[very thick, primary] (0,-1.1)--(0,-1.9);
  \fill (0,0) circle (1.8pt) node[below left] {\small $\zero$};
  \draw[->, thick, warncolor] (0,0)--(1.2,0.5) node[right] {$\vw$};
  \draw[->, thick, excolor] (0,0)--(0,1.5) node[left] {$\vz$};
\end{tikzpicture}
\end{center}
\end{example}

\begin{note}
$\vz\in W^\perp$ اگر و تنها اگر $\vz$ بر \textbf{یک مجموعه‌ی مولد} $W$ عمود باشد؛ چون آن‌وقت با خطی بودن بر هر ترکیب خطی آن‌ها هم عمود است.
\end{note}

\begin{theorem}\label{thm:perp-subspace}
$W^\perp$ زیرفضای $\R^n$ است.
\end{theorem}
\begin{proof}
$\zero\T\vw=0$، پس $\zero\in W^\perp$. اگر $\vz_1,\vz_2\in W^\perp$ و $\alpha$ اسکالر باشد، برای هر $\vw\in W$:
$(\alpha\vz_1+\vz_2)\T\vw=\alpha\,\vz_1\T\vw+\vz_2\T\vw=0$.
\end{proof}

\begin{theorem}\label{thm:perp-cap}
$W\cap W^\perp=\{\zero\}$.
\end{theorem}
\begin{proof}
اگر $\vz\in W\cap W^\perp$، $\vz$ بر هر بردار $W$ و از جمله بر خودش عمود است: $\norm{\vz}^2=\vz\T\vz=0$؛ پس $\vz=\zero$.
\end{proof}

پس طبق قضیه‌ی \ref{thm:direct-sum}، جمع $W+W^\perp$ همیشه مستقیم است.

\begin{warning}
دو زیرفضا می‌توانند متعامد باشند بدون این‌که مکمل هم باشند: در $\R^3$، $W_1=\Span\{(1,0,0)\}$ و $W_2=\Span\{(0,1,0)\}$ متعامدند، اما $W_1^\perp$ کل صفحه‌ی $yz$ است، نه فقط محور $y$.
\end{warning}

\section{الگوریتم گرام-اشمیت}

هدف: از $\va_1,\dots,\va_k$ یک \textbf{پایه‌ی یکامتعامد} برای $\Span\{\va_1,\dots,\va_k\}$ بسازیم. ایده ساده است: هر بردار جدید را برمی‌داریم، \emph{سایه‌اش} روی بردارهای ساخته‌شده‌ی قبلی را از آن کم می‌کنیم تا بر همه‌شان عمود شود، و بعد طولش را ۱ می‌کنیم.

\begin{center}
\begin{tikzpicture}[scale=0.66,>=Latex]
  \draw[gray!50] (0,0) circle (1);
  \draw[->, very thick, primary] (0,0)--(0.35,1.45) node[left] {$\va_1$};
  \draw[->, very thick, primary] (0,0)--(2.6,1.1) node[right] {$\va_2$};
  \node at (0.6,-1.8) {\small \fa{۱. بردارهای اولیه}};
  \begin{scope}[xshift=5cm]
  \draw[gray!50] (0,0) circle (1);
  \draw[->, very thick, thmcolor] (0,0)--(0.236,0.972) node[left] {$\vq_1$};
  \draw[->, very thick, primary] (0,0)--(2.6,1.1) node[right] {$\va_2$};
  \node at (0.6,-1.8) {\small \fa{۲. نرمال‌سازی $\va_1$}};
  \end{scope}
  \begin{scope}[xshift=10cm]
  \draw[gray!50] (0,0) circle (1);
  \draw[->, very thick, thmcolor] (0,0)--(0.236,0.972) node[left] {$\vq_1$};
  \draw[->, very thick, primary] (0,0)--(2.6,1.1) node[above] {$\va_2$};
  \draw[->, very thick, warncolor] (0,0)--(2.2,-0.54) node[below] {$\tilde{\vq}_2$};
  \draw[->, thick, excolor] (2.6,1.1)--(2.2,-0.54) node[midway, right] {\scriptsize $-(\vq_1\T\va_2)\vq_1$};
  \node at (0.6,-1.8) {\small \fa{۳. حذف سایه روی $\vq_1$}};
  \end{scope}
  \begin{scope}[xshift=16.6cm]
  \draw[gray!50] (0,0) circle (1);
  \draw[->, very thick, thmcolor] (0,0)--(0.236,0.972) node[left] {$\vq_1$};
  \draw[->, very thick, thmcolor] (0,0)--(0.97,-0.24) node[below right] {$\vq_2$};
  \node at (0.3,-1.8) {\small \fa{۴. نرمال‌سازی}};
  \end{scope}
\end{tikzpicture}
\end{center}

\begin{definition}[الگوریتم گرام-اشمیت]
\term{الگوریتم گرام-اشمیت}{Gram--Schmidt Algorithm}: برای $i=1,\dots,k$:
\begin{enumerate}
  \item \textbf{متعامدسازی:} $\tilde{\vq}_i=\va_i-(\vq_1\T\va_i)\vq_1-\dots-(\vq_{i-1}\T\va_i)\vq_{i-1}$.
  \item \textbf{آزمون وابستگی:} اگر $\tilde{\vq}_i=\zero$، متوقف شو.
  \item \textbf{نرمال‌سازی:} $\vq_i=\dfrac{\tilde{\vq}_i}{\norm{\tilde{\vq}_i}}$.
\end{enumerate}
\end{definition}

مثلاً سه قدم اول:
\[
\vq_1=\frac{\va_1}{\norm{\va_1}},\qquad
\tilde{\vq}_2=\va_2-(\vq_1\T\va_2)\vq_1,
\]
\[
\tilde{\vq}_3=\va_3-(\vq_1\T\va_3)\vq_1-(\vq_2\T\va_3)\vq_2 .
\]

\begin{example}
برای $\va=(1,0,1)$، $\vb=(1,0,0)$ و $\vc=(2,1,0)$ یک مجموعه‌ی متعامد (و یکامتعامد) بسازید.
\end{example}
\begin{solution}
$\vq_1=\frac{1}{\sqrt2}(1,0,1)$.

$\vq_1\T\vb=\frac1{\sqrt2}$، پس $\tilde{\vq}_2=(1,0,0)-\frac12(1,0,1)=(\frac12,0,-\frac12)$ و $\vq_2=\frac1{\sqrt2}(1,0,-1)$.

$\vq_1\T\vc=\sqrt2$ و $\vq_2\T\vc=\sqrt2$، پس
$\tilde{\vq}_3=(2,1,0)-(1,0,1)-(1,0,-1)=(0,1,0)$ و $\vq_3=(0,1,0)$.

مجموعه‌ی متعامد: $\{(1,0,1),\ (1,0,-1),\ (0,1,0)\}$؛ با تقسیم بر طول‌ها یکامتعامد می‌شود.
\end{solution}

\begin{theorem}[درستی گرام-اشمیت]\label{thm:gram-schmidt}
اگر $\va_1,\dots,\va_k$ مستقل باشند، الگوریتم متوقف نمی‌شود و برای هر $i$:
\begin{enumerate}
  \item $\vq_1,\dots,\vq_i$ یکامتعامدند؛
  \item $\vq_i$ ترکیب خطی $\va_1,\dots,\va_i$ است و $\va_i$ ترکیب خطی $\vq_1,\dots,\vq_i$؛
  \item $\Span\{\vq_1,\dots,\vq_i\}=\Span\{\va_1,\dots,\va_i\}$.
\end{enumerate}
پس $\{\vq_1,\dots,\vq_k\}$ پایه‌ای یکامتعامد برای $\Span\{\va_1,\dots,\va_k\}$ است.
\end{theorem}
\begin{proof}
با استقرا روی $i$. برای $i=1$: $\va_1\neq\zero$ (مستقل‌اند)، پس $\vq_1=\va_1/\norm{\va_1}$ طول ۱ دارد و $\Span\{\vq_1\}=\Span\{\va_1\}$.

فرض کنید حکم برای $i-1$ درست باشد.

\textbf{متوقف نمی‌شود:} اگر $\tilde{\vq}_i=\zero$ بود، آن‌گاه
\[
\va_i=\sum_{j<i}(\vq_j\T\va_i)\vq_j\in\Span\{\vq_1,\dots,\vq_{i-1}\}=\Span\{\va_1,\dots,\va_{i-1}\},
\]
که با استقلال تناقض دارد.

\textbf{تعامد:} برای $j<i$، چون $\vq_j\T\vq_l=\delta_{jl}$:
\[
\vq_j\T\tilde{\vq}_i=\vq_j\T\va_i-\sum_{l<i}(\vq_l\T\va_i)\,\vq_j\T\vq_l=\vq_j\T\va_i-\vq_j\T\va_i=0 .
\]
پس $\vq_i$ بر $\vq_1,\dots,\vq_{i-1}$ عمود است و طولش ۱ است.

\textbf{ترکیب خطی:} $\va_i=\norm{\tilde{\vq}_i}\vq_i+\sum_{j<i}(\vq_j\T\va_i)\vq_j$ ترکیبی از $\vq_1,\dots,\vq_i$ است. برعکس، $\tilde{\vq}_i$ ترکیبی از $\va_i$ و $\vq_1,\dots,\vq_{i-1}$ است و این‌ها طبق فرض استقرا ترکیبی از $\va_1,\dots,\va_{i-1}$ هستند؛ پس $\vq_i$ ترکیبی از $\va_1,\dots,\va_i$ است. از این دو، تساوی پوشش‌ها نتیجه می‌شود.
\end{proof}

\begin{note}[گرام-اشمیت به‌عنوان آزمون استقلال]
\begin{itemize}
  \item اگر الگوریتم زود متوقف نشود، $\va_1,\dots,\va_k$ مستقل‌اند (طبق قضیه‌ی \ref{thm:orth-indep}، $\vq_i$ها مستقل‌اند و همان فضا را تولید می‌کنند).
  \item اگر در قدم $j$ متوقف شود، $\va_j$ ترکیب خطی بردارهای قبلی است و بردارها وابسته‌اند:
  \[
  \va_j=(\vq_1\T\va_j)\vq_1+\dots+(\vq_{j-1}\T\va_j)\vq_{j-1}.
  \]
\end{itemize}
\end{note}

\begin{note}[پیچیدگی: $O(nk^2)$]
در قدم $i$، $i-1$ ضرب داخلی (هر کدام حدود $2n$ عمل) و $i-1$ تفریقِ مضرب یک بردار (هر کدام حدود $2n$ عمل) داریم؛ یعنی حدود $4n(i-1)$ عمل. نرمال‌سازی هم حدود $3n$ عمل است. جمع روی $i$:
$\sum_{i=1}^k 4n(i-1)\approx 2nk^2$؛ پس هزینه از مرتبه‌ی $O(nk^2)$ است.
\end{note}

\begin{corollary}[وجود پایه‌ی یکامتعامد]\label{cor:onb-exists}
هر فضای ضرب داخلی متناهی‌البُعد یک پایه‌ی یکامتعامد دارد.
\end{corollary}
\begin{proof}
فضا یک پایه‌ی متناهی دارد. گرام-اشمیت را (با $\inner{\cdot}{\cdot}$ به جای $\vq_j\T\va_i$) روی آن اجرا می‌کنیم؛ طبق قضیه‌ی \ref{thm:gram-schmidt} خروجی پایه‌ای یکامتعامد برای همان فضاست.
\end{proof}

\section{دوباره چهار زیرفضای اساسی}

در فصل ۷ دیدیم $\Nul(A)\perp\Row(A)$ و $\Nul(A\T)\perp\Col(A)$ (قضیه‌ی \ref{thm:orthogonal-subspaces}). حالا می‌توانیم دقیق‌تر بگوییم.

\begin{theorem}\label{thm:four-complements}
برای $A_{m\times n}$ حقیقی:
\begin{keyformula}
\[
\Nul(A)=\Row(A)^\perp\subseteq\R^n,
\]
\[
\Nul(A\T)=\Col(A)^\perp\subseteq\R^m .
\]
\end{keyformula}
\end{theorem}
\begin{proof}
سطرهای $A$ فضای $\Row(A)$ را تولید می‌کنند. طبق نکته‌ی بالا، $\vz\in\Row(A)^\perp$ اگر و تنها اگر $\vz$ بر هر سطر $A$ عمود باشد؛ یعنی هر درایه‌ی $A\vz$ صفر باشد: $A\vz=\zero$. بخش دوم همین استدلال برای $A\T$ است (سطرهای $A\T$ ستون‌های $A$ هستند).
\end{proof}

\begin{intuition}
هر $\vx\in\R^n$ به شکل $\vx=\vx_r+\vx_n$ با $\vx_r\in\Row(A)$ و $\vx_n\in\Nul(A)$ نوشته می‌شود و $A\vx=A\vx_r$ (شکل فصل ۷). بُعدها هم جور درمی‌آیند: $\dim\Row(A)+\dim\Nul(A)=r+(n-r)=n$.
\end{intuition}

\section{جمع‌بندی و تمرین}

\begin{summary}
\begin{itemize}
  \item یکامتعامد یعنی $\va_i\T\va_j=\delta_{ij}$. هر مجموعه‌ی متعامدِ ناصفر مستقل است.
  \item ضرایب ترکیب خطی: $c_j=\frac{\vy\T\va_j}{\va_j\T\va_j}$؛ برای یکامتعامد $c_j=\va_j\T\vy$.
  \item $W^\perp$ زیرفضاست و $W\cap W^\perp=\{\zero\}$. متعامد بودن با مکمل بودن فرق دارد.
  \item گرام-اشمیت: کم کردن سایه‌ها و نرمال‌سازی؛ پوشش را حفظ می‌کند، وابستگی را تشخیص می‌دهد و $O(nk^2)$ هزینه دارد.
  \item هر فضای ضرب داخلی متناهی‌البُعد پایه‌ی یکامتعامد دارد.
  \item $\Nul(A)=\Row(A)^\perp$ و $\Nul(A\T)=\Col(A)^\perp$.
\end{itemize}
\end{summary}

\subsection*{تمرین‌ها}
\begin{enumerate}
  \item آیا $\{(1,1,0),(1,-1,1),(-1,1,2)\}$ متعامد است؟ آن را یکامتعامد کنید.
  \item $W^\perp$ را برای $W=\Span\{(1,2,-1)\}$ در $\R^3$ بیابید.
  \item گرام-اشمیت را روی $(1,1)$ و $(2,0)$ اجرا کنید.
  \item گرام-اشمیت را روی $(1,2,3)$، $(2,4,6)$ اجرا کنید. چه اتفاقی می‌افتد؟
\end{enumerate}

\subsection*{پاسخ‌های کوتاه}
\begin{enumerate}
  \item بله؛ ضرب‌های داخلی $0$، $0$ و $0$ هستند. با تقسیم بر طول‌ها ($\sqrt2$، $\sqrt3$، $\sqrt6$) یکامتعامد می‌شود.
  \item $W^\perp=\{(x,y,z): x+2y-z=0\}$؛ یک صفحه، با پایه‌ی $\{(-2,1,0),(1,0,1)\}$.
  \item $\vq_1=\frac1{\sqrt2}(1,1)$؛ $\tilde{\vq}_2=(2,0)-(1,1)=(1,-1)$ و $\vq_2=\frac1{\sqrt2}(1,-1)$.
  \item $\tilde{\vq}_2=\zero$ و الگوریتم در قدم ۲ متوقف می‌شود؛ چون $(2,4,6)=2(1,2,3)$ و بردارها وابسته‌اند.
\end{enumerate}
